\section{Implementation}

\subsection{Architecture overview}

GeoIntel follows a classic three-tier layout: a React~\cite{react} single-page client, an Express~\cite{express} HTTP API in TypeScript, and pluggable upstream adapters that normalize heterogeneous feeds into shared schemas before the UI renders them on a CesiumJS globe~\cite{cesium}.

\begin{enumerate}[leftmargin=*]
  \item \textbf{Client (}\texttt{client/}\textbf{):} React and TypeScript; CesiumJS for the 3D map; panels for filters, events, trackers, and the Presidential Daily Brief (PDB).
  \item \textbf{Server (}\texttt{server/}\textbf{):} Express registers REST routes under \texttt{/api/...}; shared adapter logic lives under \texttt{src/adapters/} with validation (e.g., Zod) where applicable.
  \item \textbf{Data plane:} External providers (news/conflict feeds, trackers, satellite TLE sources) are fetched server-side only; credentials stay in environment variables loaded at startup (\texttt{dotenv}).
\end{enumerate}

Development uses \texttt{tsx} to run TypeScript directly; production builds bundle the server (\texttt{npm run build}, \texttt{npm start}). The client is developed with Vite and served through the same dev server in typical setups described in the repository README.

\subsection{Globe, terrain, and imagery}

The map uses \textbf{CesiumJS}. Optional \textbf{Cesium ion} terrain and OSM-style building enhancements are toggled via client env vars (e.g., \texttt{VITE\_CESIUM\_ION\_TOKEN}, \texttt{VITE\_CESIUM\_ENABLE\_TERRAIN}, \texttt{VITE\_CESIUM\_ENABLE\_BUILDINGS}). Without a token, the globe still runs with reduced 3D fidelity.

Brief sections can attach \textbf{grounded visuals}: map focus derived from modeled regions plus article imagery extracted from citation pages when embeddable images exist. That pairs narrative output with geographic and visual context.

\subsection{Feeds and trackers}

Each external source has an adapter module under \texttt{src/adapters/}. Examples include GDELT-style document queries~\cite{gdelt}, ACLED where credentials allow~\cite{acled}, AIS streams and AISHub for maritime context, OpenSky and optional supplemental aircraft sources~\cite{opensky}, and CelesTrak-based TLE propagation for satellites~\cite{celestrak}. Responses are cached server-side with TTLs to respect rate limits and stability.

\subsection{Generative AI: Presidential Daily Brief}

The PDB workflow calls Google's \textbf{Gemini} models through the Gemini API~\cite{gemini-api}. When search grounding is enabled, the server sends requests that include the Google Search tool so outputs can cite retrieved URLs.
The client renders citations as links and persists successful section payloads in browser storage for resilience across reloads.

Generation proceeds \textbf{section by section} (one request per thematic slide), which is a deliberate multi-call pipeline rather than a single monolithic completion: it trades extra latency for clearer failure isolation and UI progress feedback.

\textbf{Agentic tooling:} We do not rely on open-ended autonomous agents browsing the web on their own. The closest analogue is grounded retrieval plus constrained JSON-oriented prompting inside fixed API routes---deterministic orchestration from the server, not emergent agent loops.

\subsection{PDF export}

Brief PDF export is handled \textbf{in the browser}: the UI captures rendered slide content and emits a PDF (client-side libraries), avoiding server-side storage of generated documents.

\subsection{Optional SpaceMouse bridge}

For 3D navigation hardware, an optional \textbf{Python} bridge (\texttt{script/spacemouse\_bridge.py}) can expose SpaceMouse motion over a local WebSocket (\texttt{127.0.0.1:8765}). \texttt{npm run dev} attempts to start this bridge when the port is free; \texttt{npm run dev:no-spacemouse} skips it.

\subsection{Configuration and secrets}

Secrets (Gemini keys, tracker keys, ion tokens, etc.) live only in \texttt{.env} or equivalent host env---never committed to reports or screenshots at full resolution if pasted verbatim. The repository ships \texttt{.env.example} as a template; operators copy it to \texttt{.env} locally.

\subsection{Reproducibility}

This write-up mirrors the deliverable expectations for the course final project (written report + code repository with runnable instructions)~\cite{course-req}.

\textbf{Repository:} \url{https://github.com/Ben-W-Yal/Good_Vibes_Final_Project}

\textbf{Steps to run locally:}
\begin{enumerate}[leftmargin=*]
  \item Clone the repo and install dependencies: \texttt{npm install}.
  \item Copy \texttt{.env.example} to \texttt{.env} and add at minimum a valid Gemini API key plus any optional provider keys you need for trackers or enhanced terrain.
  \item Start development: \texttt{npm run dev} (or \texttt{npm run dev:no-spacemouse} without the SpaceMouse bridge).
  \item Open the served URL (default port shown in server logs, commonly \texttt{localhost:3000} in development).
\end{enumerate}
